\documentclass{article}

\usepackage[utf8]{inputenc}
\usepackage{graphicx}
\usepackage{url}
\usepackage{geometry}
\usepackage{setspace} % for spacing control
\usepackage{parskip}  % better paragraph spacing

% Margins to match your original letter margins
\geometry{
    top=1in,
    bottom=1in,
    left=1.5in,
    right=1.5in
}

\bibliographystyle{plain}

\begin{document}


% Logo top left
\begin{flushleft}
    \includegraphics[width=2.8in]{figures/ifisc-marcas-horizontal.pdf}
\end{flushleft}

\vspace{0.1cm}


\vspace{1cm}

% Date (optional, uncomment if needed)
%\begin{flushright}
%    \today
%\end{flushright}

% Greeting
Dear Editor,

\vspace{0.5cm}

% Body of the letter
Please find together with this letter our manuscript entitled ``Mobility energetic balance frames spatial configuration of cities" for your consideration as an Article in Nature Cities.

Mobility plays a central role in everyday activity in cities. It is essential not only for the functioning of work places, business and services, but for the citizens who require such services as well. In this context,  mobility has been used as an input when studying the well-being and livability of urban areas, from sprawl indices to accessibility to services, segregation of ethnic minorities to the emissions and air quality. 

Here we maintain a global view over city scale, but use mobility in a new way, as an indicator of the energetic costs that an urban system must sustain. These costs refer to personal effort when biking or walking, or other metrics such as gas or electricity consumption if we consider mobility in private vehicles or public transportation. The common feature to them all is that it is a price to pay for each day of activity of a city. We take an approximation and calculate the mobility energetic costs of full cities, 57 worldwide, and show that the spatial configuration of the urban area is such that minimizes energetic costs. This is a general rule in all the cities observed except for a few exceptions where there are neighboring areas not accessible for expansion due to wetlands and other geographical constraints. This general rule allow us to propose a method based on mobility energetic costs in the planning of future urban expansions. 

All in all, the relevance of our general finding, its implication for the new science of cities and its applicability in urban planning make us believe that our work not only fall within the scope of Nature Cities but it is of interest for your multidisciplinary audience.

For your convenience, we include here a list of potential reviewers:

\begin{itemize}
    \item \textbf{Prof. Marta González} \\
    Affiliation: University of California, Berkeley (UC Berkeley), USA. \\
    E-mail: \texttt{martag@berkeley.edu}
    
    \item \textbf{Prof. Marc Barth\'elemy} \\
    Affiliation: CEA Saclay, Paris, France. \\
    E-mail: \texttt{marc.barthelemy@cea.fr}
    
    \item \textbf{Prof. Rafael Prieto-Curiel} \\
    Affiliation: Complexity Science HUB, Vienna, Austria. \\
    E-mail: \texttt{ prieto-curiel@csh.ac.at }
    
    \item \textbf{Prof. Geoff Boeing} \\
    Affiliation: University of Southern California, Los Angeles, USA. \\
    E-mail: \texttt{boeing@usc.edu}
    
    \item \textbf{Prof. Elsa Arcaute} \\
    Affiliation: Centre for Advanced Spatial Analysis, University College London, UK. \\
    E-mail: \texttt{e.arcaute@ucl.ac.uk}
\end{itemize}
 \smallskip
Thank you for considering our submission. We look forward to your response. 

\vspace{0.5cm}

Yours sincerely,\\

Federico Bellisardi, Enrique Cano, Ana Ruiz-Varona and Jos\'e J. Ramasco

\vspace{1cm}

\end{document}
